\documentclass[11pt]{article}
\usepackage{ochamath}
\subjectcolor{2F5DA8}
\setsheet{線形代数}{二次形式と正定値　演習}{https://university-math-crj.pages.dev/linear-algebra/quadratic-forms/}
\begin{document}
\makesheethead{解答}

\problem[展開と平方完成]
\begin{ans}
$q=3x^2+2xy+2y^2=3(x+y/3)^2+(5/3)y^2$。非零 $(x,y)$ では少なくとも一方の2乗が正なので正定値。左上3と行列式5が正であることでも確認できる。
\end{ans}
\point{交差項は対称行列の非対角成分の2倍。}

\problem[主軸変換]
\begin{ans}
固有値7と3に対して、固有ベクトルは $(1,1)^{\mathsf T}/\sqrt2$ と $(1,-1)^{\mathsf T}/\sqrt2$。$z_1=(x+y)/\sqrt2,\ z_2=(x-y)/\sqrt2$ により $q=7z_1^2+3z_2^2$。展開すると $5x^2+4xy+5y^2$ に戻る。
\end{ans}
\point{固有値が新しい座標の2乗の係数になる。}
\newpage

\problem[符号の分類]
\begin{ans}
前者は $4x^2$ で半正定値だが、非零ベクトル $(0,1)$ でも0なので正定値ではない。後者は $2x^2-3y^2$ で、$(1,0)$ では2、$(0,1)$ では $-3$ となり、不定値。
\end{ans}
\point{非零なのに値が0になる方向も調べる。}

\problem[ばねのエネルギー]
\begin{ans}
$U=\frac12\times2\times(0.01)^2=0.0001\ \mathrm J$。$K$ の固有値は1,3でともに正なので、非零変位のエネルギーは正。単位は $\mathrm{m}\times\mathrm{N/m}\times\mathrm m=\mathrm J$。
\end{ans}
\point{エネルギーの式の1/2を忘れない。}
\end{document}
